
\section{Appendix B. Implementation Details}

\subsection{B.1 The `NystromAttention` Class}

\begin{lstlisting}
class NystromAttention(BaseAttention):
    def \_\_init\_\_(self, config=None, rank=None, num\_landmarks=None,
                 pinv\_iterations=3, pretrained\_attn=None, old\_attn=None,
                 comp\_config=None, **kwargs):
        # Dimensions
        self.dim = ...
        self.num\_heads = ...
        self.num\_kv\_heads = ...
        self.head\_dim = ...
        self.all\_head\_size = ...
        self.all\_kv\_head\_size = ...
        self.kv\_repeat = ...
        
        # Landmarks
        self.num\_landmarks = min(32, int(num\_landmarks or 32))
        self.pinv\_iterations = pinv\_iterations
        
        # Decoder detection
        self.is\_decoder = ...
        self.is\_nested\_encoder = False
        
        # LoRA config
        self.lora\_rank = ...
        
        # Projections
        self.query = LoRALinear(...)
        self.key = LoRALinear(...)
        self.value = LoRALinear(...)
        self.out\_proj = LoRALinear(...)
        
        if source\_attn is not None:
            self.copy\_weights(source\_attn)
\end{lstlisting}

\subsection{B.2 The Iterative Pseudoinverse}

\begin{lstlisting}
def \_iterative\_pinv(self, A, iters=3):
    try:
        return torch.linalg.pinv(A.float()).to(A.dtype)
    except Exception:
        norm\_1 = A.abs().sum(dim=-2).amax(dim=-1, keepdim=True).unsqueeze(-1)
        norm\_inf = A.abs().sum(dim=-1).amax(dim=-1, keepdim=True).unsqueeze(-1)
        alpha = 1.0 / (norm\_1 * norm\_inf + 1e-5)
        
        V = alpha * A.transpose(-2, -1)
        I = torch.eye(A.size(-1), device=A.device, dtype=A.dtype).unsqueeze(0).unsqueeze(0)
        
        for \_ in range(iters):
            V = torch.matmul(V, 2.0 * I - torch.matmul(A, V))
        return V
\end{lstlisting}

\textbf{Note.} We first try \texttt{torch.linalg.pinv} (SVD). If it fails, we fall back 
to Newton-Schulz. This is a deviation from the original paper.

\subsection{B.3 The Three Execution Paths}

\begin{lstlisting}
def forward(self, hidden\_states, attention\_mask=None, ...):
    Q = self.query(hidden\_states).view(B, N, H, d).transpose(1, 2)
    K = self.key(hidden\_states).view(B, N, Hkv, d).transpose(1, 2)
    V = self.value(hidden\_states).view(B, N, Hkv, d).transpose(1, 2)
    
    # RoPE
    if position\_embeddings is not None:
        Q, K = apply\_rotary\_pos\_emb(Q, K, cos, sin)
    
    # GQA
    if self.kv\_repeat \textgreater{} 1:
        K = repeat\_kv(K, self.kv\_repeat)
        V = repeat\_kv(V, self.kv\_repeat)
    
    valid\_mask = get\_token\_valid\_mask(attention\_mask, B, N, device)
    
    # PATH 1: Causal decoder
    if self.is\_decoder or self.is\_causal:
        # Dense causal
        ...
        return out, None
    
    # PATH 2: Short sequence
    if N \textless{}= max(128, 2 * self.num\_landmarks):
        context = self.\_dense\_attention(Q, K, V, valid\_mask, dtype)
        return out, None
    
    # PATH 3: Nystr m
    # adaptive pooling landmarks
    Q\_landmarks = F.adaptive\_avg\_pool1d(Q\_flat, m)...
    K\_landmarks = F.adaptive\_avg\_pool1d(K\_flat, m)...
    
    scores1 = Q @ K\_landmarks^T
    scores2 = Q\_landmarks @ K\_landmarks^T
    scores3 = Q\_landmarks @ K^T
    
    F1 = softmax(scores1); F2 = softmax(scores2); F3 = softmax(scores3)
    F2\_pinv = self.\_iterative\_pinv(F2)
    
    step1 = F3 @ V
    step2 = F2\_pinv @ step1
    context = F1 @ step2
    
    # Head mask, valid mask, out\_proj
    ...
\end{lstlisting}

\subsection{B.4 Landmark Selection via Adaptive Pooling}

\begin{lstlisting}
m = min(self.num\_landmarks, N)

Q\_flat = Q.reshape(B * H, N, d).transpose(1, 2)   # (BH, d, N)
K\_flat = K.reshape(B * H, N, d).transpose(1, 2)

Q\_landmarks = F.adaptive\_avg\_pool1d(Q\_flat, m).transpose(1, 2).reshape(B, H, m, d)
K\_landmarks = F.adaptive\_avg\_pool1d(K\_flat, m).transpose(1, 2).reshape(B, H, m, d)
\end{lstlisting}

\textbf{Why adaptive\_avg\_pool1d?} It partitions the sequence into $m$ segments 
of approximately equal length and averages each. This is deterministic and 
preserves locality.

\subsection{B.5 Weight Transfer}

Four patterns:

\begin{enumerate}
\item \textbf{GPT-2} -- split fused \texttt{c\_attn}, \texttt{c\_proj} for output
\item \textbf{RoBERTa} -- \texttt{self.query/key/value}, \texttt{output.dense}
\item \textbf{Qwen/DistilBERT} -- flat projections
\item \textbf{UniversalAdapter} -- fallback

\end{enumerate}
\subsection{B.6 Injection Hook}

\begin{lstlisting}
def inject\_nystrom(layer, model\_config, comp\_config, old\_attn=None, **kwargs):
    src\_attn = old\_attn or layer.attention...
    hidden\_size = ...
    num\_heads = ...
    
    try:
        landmarks = getattr(comp\_config, "num\_landmarks", 32) if comp\_config else 32
        new\_attn = NystromAttention(
            config=model\_config,
            num\_landmarks=landmarks,
            old\_attn=src\_attn,
            hidden\_size=hidden\_size,
            num\_heads=num\_heads,
        )
        if src\_attn is not None:
            new\_attn.copy\_weights(src\_attn)
        
        if hasattr(layer, "attention"):
            if hasattr(layer.attention, "self"):
                layer.attention.self = new\_attn
                new\_attn.is\_nested\_encoder = True
            else:
                layer.attention = new\_attn
        elif hasattr(layer, "self"):
            layer.self = new\_attn
        elif hasattr(layer, "attn"):
            layer.attn = new\_attn
        return new\_attn
    except Exception as e:
        logger.error(f"injection failed: {e}")
        return None
\end{lstlisting}

\subsection{B.7 Common Pitfalls}

\subsubsection{B.7.1 Causal Without Fallback}

Using Nystr m for causal models leaks future information. Must use 
dense causal.

\subsubsection{B.7.2 No Short Sequence Fallback}

For $N \leq 2m$, dense is faster and more accurate.

\subsubsection{B.7.3 Conv1D Not Transposed}

GPT-2's Conv1D stores \texttt{(in, out)}; must transpose.

\subsubsection{B.7.4 Pinv Without Preconditioning}

Newton-Schulz requires $\|\alpha A\| < 1$. Without preconditioning, it 
diverges.

\subsubsection{B.7.5 Wrong Associative Order}

Computing $F_1 \cdot \text{pinv}(F_2)$ before $F_3 V$ gives $\mathcal{O}(N^2 m)$, 
defeating the purpose.

\subsection{B.8 Numerical Stability}

\begin{itemize}
\item Pinv computed in FP32.
\item \texttt{torch.nan\_to\_num} guards.
\item Landmark pooling in input dtype.

\end{itemize}
\subsection{B.9 Comparison with Companion Methods}

\begin{lstlisting}

---
\end{lstlisting}